\chapter*{Conclusion}
\addcontentsline{toc}{chapter}{Conclusion}
\markboth{Conclusion}{Conclusion}

The goal of this thesis was to build on the bachelor-thesis prototype of an autonomous agent for \Trackmania{} and create a more systematically trained and evaluated version. We kept the main idea that the agent does not have to perceive the game through images, but can instead work with the car state, reconstructed track geometry, and virtual sensors. In the master's thesis, we extended this idea into a complete cycle: obtaining state from the game, computing a compact observation, applying an action through a virtual controller, and automatically training the behavior policy.

The experiments in the tested conditions showed that the proposed observation is usable for a neural network. Supervised learning helped select a practical architecture and indicated that the network can imitate human actions, but it also revealed the limits of behavioral cloning alone. Neuroevolution therefore became the main training path. In our experiments, the most readable evaluation of driving behavior was a lexicographic ordering of goals: first finish the track, then get as far as possible, then improve time, and only afterwards distinguish the number of crashes. This approach avoided an opaque weighted sum of constants and provided a more stable signal than the tested alternatives.

The final agent improved substantially over the bachelor-thesis version. We evaluated a selected individual from training on \texttt{single\_surface\_flat} on the unseen map \texttt{small\_map} without additional fine-tuning, where it achieved \(19.68\,\mathrm{s}\). The bachelor-thesis agent achieved \(23.064\,\mathrm{s}\) on the same map, while the human reference run had \(17.89\,\mathrm{s}\). On the training map \texttt{single\_surface\_flat}, the best agent after longer training reached \(25.92\,\mathrm{s}\), close to the mean of the saved player-time sample. The agent therefore still does not outperform the best human players, but in the tested conditions it successfully approached typical player performance and clearly surpassed the original prototype. The final training runs also showed that the method is not tied only to one flat track: finishing policies were found both on a track with height differences and on a track with different surfaces.

\section*{Fulfillment of Goals and Contributions}

The main goals of the assignment can be considered fulfilled. We designed and implemented a new version of the agent that uses a geometric environment representation, automated genetic training, and systematic logging of results. The agent improved not only in time, but also in the fact that its behavior can be analyzed using metrics, plots, and trajectories.

The contribution of the thesis is not only a single best time. The important result is the whole solution: turning a digital track into an observable environment, deriving measurable criteria from driving, and selecting from a population of candidates a policy that can finish the track autonomously. The thesis also showed which directions were useful and which had limitations: supervised learning was useful for initialization and architecture selection, interactive imitation improved the data, but autonomous improvement of the agent was best supported by lexicographic neuroevolution.

\section*{Limitations and Future Work}

The results should be interpreted cautiously. Most experiments are practical comparisons with one random setting, so we do not draw strong statistical conclusions about the general optimality of hyperparameters. The main reason is computational cost: some experiments, for example the search for a suitable mutation probability and mutation variance, took hundreds of hours of agent simulation in the environment. Repeating the same grid experiments several times would be very demanding within the resources of a master's thesis. The agent was also tested on a selected set of maps, not on a large collection of tracks. The geometric observation also depends on the quality of the reconstructed map and on which properties of the world are included in the input.

Further development should focus on broader evaluation on more maps and independent runs. It would also be interesting to examine hybrid training more deeply, where human data provide a good initial direction and neuroevolution then optimizes behavior directly according to driving performance. A more detailed comparison with reinforcement learning would also make sense, since PPO showed a usable signal in this thesis, although it was not the main tuned direction.

The thesis showed that an autonomous agent for \Trackmania{} does not have to be built only on image input or on manually designed rules. If it has access to a meaningful geometric representation of the world and a training procedure that evaluates complete drives, it can learn behavior usable in a real game environment.
